% SOURCE_AUTHORITY: sga5_fr_workpass.tex, lines 9292--9320 (LF-normalized unit).
% SOURCE_SHA256: 791F4EFFC5E02832D5D77ED03518C8156D6F07E4C8238B03545DB93D883FBB28
\begin{statement}{Proposição 2.2.2}
Sob as hipóteses anteriores, temos:
\begin{enumerate}[label=\alph*)]
\item
\[
\R^{2i+1}p_*(A_P)=0\quad\text{Para tudo }i\in\mathbb N.
\]
\item Quando $\bigoplus_i \R^{2i}p_*(\mu_P^{\otimes i})$ é dotado de sua estrutura de $A_S$-álgebra derivada do produto cup, o morfismo de $A_S$-álgebra
\[
A_S[T]\longrightarrow \bigoplus_i \R^{2i}p_*(\mu_P^{\otimes i})
\]
definido por $\eta$ tem por núcleo o ideal $(T^{r+1})$ e define, portanto, um isomorfismo de $A_S$-álgebra
\[
\theta:A_S[T]/(T^{r+1})\xrightarrow{\sim}\bigoplus_i \R^{2i}p_*(\mu_P^{\otimes i}).
\]
\item O morfismo
\[
A_S\longrightarrow \R^{2r}p_*(\mu^{\otimes r})
\]
Obtido compondo $\theta$ com a restrição a $A_S$ da multiplicação por $T^r$ de $A_S[T]/(T^{r+1})$ em si é o inverso do morfismo de traço (SGA 4 XVIII)
\[
\R^{2r}p_*(\mu^{\otimes r})\longrightarrow A_S;
\]
Em outras palavras, temos
\[
\operatorname{Tr}(\eta^r)=1.
\]
\end{enumerate}
\end{statement}
